
\subsubsection{4.1 Existence Gate Experiment (h-e1)}

The first causal mechanism step is tested: do static analyzers produce actionable warnings on benchmark code?

\textbf{Configuration:}
\begin{itemize}
\item Dataset: HumanEval (164 problems, full test set)
\item Code source: Canonical solutions (proxy for LLM output due to API unavailability)
\item Analyzer: pylint with \texttt{--disable=C,R}
\item Gate threshold: $\geq 30$\% of problems with $\geq 1$ warning
\end{itemize}

\textbf{Proxy rationale:} Without OpenAI/Claude API access, canonical solutions were used as a lower-bound proxy. Canonical solutions are well-written human code; LLM-generated code typically exhibits more issues, so this represents a floor for expected warning rates.

\subsubsection{4.2 Implementation}

\textbf{Pipeline components:}
\begin{itemize}
\item HumanEval loader: Loads problems from HuggingFace (\texttt{openai\textbackslash \_humaneval})
\item Static analyzer: Runs pylint, parses JSON output, filters by severity
\item Metrics aggregator: Computes warning counts and rates
\end{itemize}

\textbf{Artifacts generated:}
\begin{itemize}
\item \texttt{metrics.json}: Aggregated warning statistics
\item \texttt{pylint\textbackslash \_results.json}: Per-problem pylint output
\item \texttt{warning\textbackslash \_distribution.png}: Histogram of warnings per problem
\item \texttt{warning\textbackslash \_type\textbackslash \_breakdown.png}: Distribution by warning category
\item \texttt{gate\textbackslash \_metrics.png}: Threshold vs. actual comparison
\item \texttt{top\textbackslash \_warning\textbackslash \_codes.png}: Most frequent warning codes
\end{itemize}
